
The findings of this study provided a critical, data-driven framework for re-evaluating the solid waste management strategies of the Municipality of Bacolod. By transitioning from static administrative heuristics to dynamic, AI-optimized policy design, the results offered specific, actionable insights that directly addressed the municipality's chronic compliance deficits. This section discussed the practical implications of these findings.

For the Municipality of Bacolod, the most significant implication of this research was the invalidation of the ``Equal Distribution'' model. The current practice of spreading the limited annual budget of \textpeso 1,500,000 thinly across all seven barangays was mathematically proven to be the primary driver of failure in high-density areas. The simulation demonstrated that this approach diluted enforcement intensity below the critical threshold required to alter behavior in urban centers like Barangay Poblacion. To address this, the LGU had to adopt the ``Sequential Saturation'' strategy discovered by the DRL agent. This required a political shift: rather than funding all barangays inadequately at the same time, the LGU should concentrate resources on one or two targets at a time—effectively ``besieging'' non-compliant zones until social norms take over—before rotating funds to the next priority area. This approach respected the strict budgetary ceiling while maximizing the per-peso impact of every intervention.

Furthermore, the study highlighted the necessity of decoupling policy instruments. The finding that incentives failed in large populations due to budget dilution suggested that the LGU should restrict monetary rewards to smaller, manageable puroks or barangays where the per-capita value remained meaningful. Conversely, in high-density zones, the budget should be redirected toward high-visibility enforcement to overcome anonymity. This targeted differentiation allowed the municipality to move away from a ``one-size-fits-all'' ordinance toward a precise, demographic-specific application of R.A. 9003.
